\begin{table}[t]
\centering
\caption{Heterogeneity in the Effect of Flood-Induced Migration}
\label{tab:heterogeneity}
\begin{tabular}{lcc}
\toprule
 & $\Delta$OccScore (1940) & $N$ \\
\midrule
\multicolumn{3}{l}{\textit{Panel A: By Age at Displacement}} \\
Age 15-25 & 0.862 & 17,517 \\
 & (2.868) & \\
Age 26-35 & 2.024 & 14,709 \\
 & (2.370) & \\
Age 36-45 & 4.033 & 8,231 \\
 & (3.799) & \\
Age 46-60 & 5.857 & 3,734 \\
 & (7.968) & \\
\midrule
\multicolumn{3}{l}{\textit{Panel B: By Pre-Flood Occupational Status}} \\
Low pre-flood status & 2.792 & 24,791 \\
 & (2.480) & \\
High pre-flood status & 1.122 & 19,400 \\
 & (2.646) & \\
\bottomrule
\end{tabular}
\par\vspace{0.3em}
{\footnotesize
\textit{Notes:} Each row reports the IV coefficient on migration (instrumented by flood exposure) from a separate 2SLS regression on the indicated subsample. The dependent variable is the change in occupational earnings score between 1920 and 1940. Panel A splits the sample by age in 1920 to test whether younger workers benefited more from flood-induced displacement. Panel B splits at the median pre-flood occupational score to test whether higher-status workers had more transferable skills. All specifications include age fixed effects and pre-flood controls. Standard errors clustered by 1920 county of residence.
}
\end{table}
